% State of the Art section. Paste into your own thesis, or \input it.
%
% Built from docs_thesis/literature.md Axis 1 and Axis 3.
% Verification status: all entries below are [BIBLIOGRAPHY VERIFIED], abstract-level.

\section{Getting the tree}
\label{sec:getting-tree}

Every feature in this thesis is computed on a tree built from a segmentation, so the segmentation
has to support that construction before any feature definition matters.
Voxel overlap with a reference mask is not the property that decides whether it does.

CAS-Net, the segmentation method this thesis builds on (Section~\ref{sec:segmentation}),
reaches a Dice similarity coefficient of 91.2 against expert annotation, close to the 92.8 two
independent annotators reach against each other \cite{bransby2026imagecasx}.
A high overlap score does not guarantee a correct tree, because Dice is computed voxel by voxel and
cannot see whether those voxels stay connected.
Shit et al.\ made that gap precise: two segmentations can share a Dice score and disagree completely
on topology, since Dice has no term for connectivity \cite{shit2021cldice}.
Their soft-clDice measure and the Betti number error it motivates, both reported in the ImageCAS-X
benchmark alongside Dice \cite{bransby2026imagecasx}, are what actually detect a break in the vessel
tree that overlap alone would miss.

A broken prediction is not merely a lower score in a pipeline built on the tree rather than the mask.
It is a structurally wrong tree: an extra connected component, a spurious terminus where a vessel was
cut, or a missing bifurcation where two components should have joined --- and every one of those
corrupts a branch count or a bifurcation angle computed downstream.
Qiu et al.\ target exactly this failure mode with a three-stage method: a centerline-enhanced
segmentation loss, a regularized-walk reconnection of broken segments using combined distance and
direction similarity, and implicit-neural reconstruction of the vessels that remain missing after
that \cite{qiu2025topology}.
On two public datasets it reaches a Dice of 88.53 and a Hausdorff distance of 1.07~mm on ASOCA, and
85.07 and 1.63~mm on PDSCA.
Those numbers are not directly comparable to the CAS-Net baseline this thesis uses, since neither
dataset is ImageCAS, so adopting the method would mean re-benchmarking it on the 160 ImageCAS-X test
cases rather than taking the published figures on faith.

What this leaves for objective 6 is a specific instruction rather than a general one: reliability
against the ground truth has to be checked for connectivity, not only for overlap, because a
segmentation can score well on Dice and still hand the topology pipeline the wrong tree.
